\input{preamble}

\begin{document}

\title{Genealogía familiar}
\maketitle

\phantomsection
\label{section-phantom}

\tableofcontents

\section{Mapa general}
\label{section-genealogia-mapa-general}

Estas notas reúnen lo que voy averiguando sobre
las familias González Casanova, Henríquez Ureña
y Lombardo Toledano. La regla aquí es distinguir
entre datos documentados, recuerdos familiares y
conjeturas pendientes. En particular, un apellido
no basta para decidir un origen geográfico.

La parte del árbol que por ahora está mejor
establecida puede visualizarse así:

\[
\xymatrix@C=1.6em@R=1.7em{
&& Pedro Henríquez Ureña
  \ar@{-}[d]
&& Isabel Lombardo Toledano
  \ar@{-}[dll]
\\
&& \multicolumn{3}{c}{\text{Natacha Henríquez Lombardo}}
  \ar@{-}[d]
\\
&& \multicolumn{3}{c}{\text{Pablo González Casanova (1922--2023)}}
}
\]

Para entender de dónde vienen esos apellidos hay
que abrir el árbol hacia arriba. Dos ramas son
especialmente claras: Henríquez Ureña conduce a
República Dominicana y Curazao, y Lombardo conduce
a un inmigrante italiano del siglo XIX. La rama
Casanova está documentada en Yucatán desde el
siglo XIX, pero su origen anterior todavía está
abierto.

\section{Los González Casanova}
\label{section-genealogia-gonzalez-casanova}

Pablo González Casanova (1922--2023), rector de
la UNAM entre 1970 y 1972, fue hijo del lingüista
y filólogo Pablo González Casanova y de Concepción
del Valle Romo. Este parentesco aparece también en
las semblanzas institucionales de la UNAM.\footnote{%
UNAM, ``Valioso legado de Don Pablo'',
\url{https://www.cch.unam.mx/comunidad/valioso-legado-de-don-pablo}.}

En la generación anterior aparece la unión de dos
líneas: González y Casanova. El padre del lingüista
Pablo González Casanova fue Pablo González Aznar y
su madre Encarnación Casanova Castillo. Por ahora,
la rama Casanova que he podido reconstruir de modo
provisional es

\[
\xymatrix@C=1.3em@R=1.5em{
José García\hbox{-}Casanova \ar@{-}[r]
  & Rosa Vázquez
\\
& Francisco García\hbox{-}Casanova Vázquez \ar@{-}[d]
\\
& Francisco García\hbox{-}Casanova Zaldívar \ar@{-}[d]
\\
& Encarnación Casanova Castillo \ar@{-}[d]
\\
& Pablo González Casanova\ (1889--1936) \ar@{-}[d]
\\
& Pablo González Casanova\ (1922--2023).
}
\]

La investigación pendiente importante es encontrar
el origen de José García-Casanova. El apellido
\emph{Casanova} existe tanto en Italia como en la
Península Ibérica; por eso no hay todavía base para
decir que esta rama sea italiana. Lo que sí sabemos
es que la familia García-Casanova estaba ya asentada
en Yucatán en el siglo XIX.

\section{Pedro Henríquez Ureña}
\label{section-genealogia-henriquez-urena}

Pedro Henríquez Ureña (1884--1946) nació en Santo
Domingo. Fue hijo de Francisco Henríquez y Carvajal
y de la poeta Salomé Ureña. Se casó con Isabel
Lombardo Toledano en México. Tuvieron dos hijas:
Natacha, nacida en México, y Sonia, nacida en
Argentina. Este último dato está confirmado por una
entrevista con Sonia Henríquez Ureña conservada por
El Colegio de México.\footnote{%
Mary Carmen Sánchez Ambriz, ``El acervo de Pedro
Henríquez Ureña a El Colegio de México. Entrevista
con Sonia Henríquez Ureña de Hlito'',
\url{https://libros.colmex.mx/wp-content/uploads/2021/04/boled_122.pdf}.}

La línea paterna de Pedro conduce a Curazao. Su
abuelo Noel Henríquez pertenecía a una familia
curazoleña. Hay una tradición familiar y
bibliográfica que relaciona esta rama con la
diáspora sefardí del Caribe y de los Países Bajos.
Éste es un punto que merece una investigación
archivística separada: antes de prolongar el árbol
hasta Europa quiero verificar los enlaces uno por
uno en registros de Curazao y, si es posible, en
archivos de comunidades sefardíes neerlandesas.

La línea materna de Pedro es la de Salomé Ureña,
una de las figuras fundamentales de la literatura y
la educación dominicanas. De este modo, esta rama
del árbol no es sólo genealógica: está directamente
ligada a la historia intelectual de la República
Dominicana y, por Pedro, a la historia intelectual
de México.

\section{Lombardo Toledano: la rama italiana}
\label{section-genealogia-lombardo-toledano}

Aquí sí aparece una procedencia italiana
explícitamente documentada. Vicente Lombardo
Toledano (1894--1968) fue hijo de Vicente Lombardo
Carpio e Isabel Toledano. Según la cronología
biográfica preparada por Rosa María Otero y Gama de
Lombardo, sus abuelos paternos fueron \emph{Vincenzo
Lombardo Cati}, originario de Turín, Italia, y
Marcelina Carpio. Vincenzo llegó a México en
1859.\footnote{%
Centro de Estudios Filosóficos, Políticos y Sociales
Vicente Lombardo Toledano, \emph{Escritos
autobiográficos}, cronología de Rosa María Otero y
Gama de Lombardo,
\url{https://www.centrolombardo.edu.mx/wp-content/uploads/formidable/14/Escritos_autobiogr%C3%A1ficos_REV.pdf}.}

La misma fuente dice que los abuelos maternos de
Vicente eran Miguel Toledano y Flora Toledano,
descendientes de una antigua familia española cuyo
fundador había llegado a México en el siglo XVII.
Esto es especialmente interesante porque impide
identificar sin más ``Lombardo Toledano'' con una
sola procedencia: la rama Lombardo es italiana,
mientras que la rama Toledano está descrita allí
como española.

Isabel Lombardo Toledano, hermana de Vicente, se
casó con Pedro Henríquez Ureña. Por tanto,

\[
\xymatrix@C=1.1em@R=1.4em{
Vincenzo Lombardo Cati \ar@{-}[r]
  & Marcelina Carpio
\\
& Vicente Lombardo Carpio \ar@{-}[d]
\\
& \begin{array}{c}
   \text{Vicente Lombardo Toledano}\quad
   \text{Isabel Lombardo Toledano}
  \end{array}
\\
&& Isabel Lombardo Toledano \ar@{-}[d]
\\
&& Natacha Henríquez Lombardo.
}
\]

El siguiente paso para esta rama es localizar el
acta de nacimiento de Vincenzo en el Piamonte y
reconstruir sus padres y abuelos. Sólo entonces se
podrá precisar más que ``originario de Turín''.

\section{Una red intelectual mexicana}
\label{section-genealogia-red-intelectual}

La genealogía se cruza de manera llamativa con la
historia intelectual mexicana. Pedro Henríquez Ureña
formó parte del núcleo del Ateneo de la Juventud
junto con figuras como Antonio Caso, Alfonso Reyes
y José Vasconcelos. Décadas después, Vicente
Lombardo Toledano fue alumno de Antonio Caso y se
convirtió en una figura central del debate político
y universitario mexicano.

Las relaciones entre Caso, Lombardo, Vasconcelos y
Henríquez Ureña fueron además personales e
institucionales, no sólo bibliográficas. La propia
cronología histórica de la UNAM registra, por
ejemplo, los conflictos universitarios de 1923
entre Lombardo y Vasconcelos, la renuncia de
Antonio Caso y la salida de Pedro Henríquez Ureña
de sus cargos docentes.\footnote{%
UNAM, cronología histórica, ``1920'', sección sobre
1923,
\url{https://www.unam.mx/acerca-de-la-unam/unam-en-el-tiempo/cronologia-historica-de-la-unam/1920}.}

Esto sugiere que una historia familiar completa no
debe ser sólo una lista de ascendientes. Hay también
una historia de redes: parentesco, amistad,
magisterio, instituciones y disputas intelectuales.

\section{Recuerdos familiares por verificar}
\label{section-genealogia-recuerdos-familiares}

Estas notas no se toman todavía como hechos
históricos. Se conservan porque pueden orientar la
búsqueda documental.

\begin{itemize}
\item En recuerdos familiares aparecen los nombres
Charles y Mara, Henrique y Lolita, Manuel y Helena,
y el apodo ``Titopaul'' para Pablo. Falta fijar con
precisión qué generación y qué parentesco designa
cada referencia.
\item Aparece también una ``tía Sonia''. Esto es
compatible con Sonia Henríquez Ureña, hermana de
Natacha, pero no debe identificarse automáticamente
sin contexto adicional.
\item Hay un recuerdo de una conexión familiar de
los Bracho--Carpizo--Valle con Antonio Caso. Todavía
no hay identificación documental suficiente para
dibujar esa rama.
\end{itemize}

\section{Preguntas abiertas}
\label{section-genealogia-preguntas-abiertas}

Las tres búsquedas principales que quedan abiertas
son:

\begin{enumerate}
\item \textbf{García-Casanova.} Encontrar el origen
de José García-Casanova y decidir documentalmente
si la línea anterior conduce a España, Italia u
otro lugar.
\item \textbf{Henríquez.} Reconstruir la línea de
Curazao generación por generación y comprobar la
tradición de ascendencia sefardí con registros
primarios.
\item \textbf{Lombardo.} Localizar a Vincenzo
Lombardo Cati en registros piamonteses y extender
la rama italiana hacia atrás.
\end{enumerate}

Estas notas deben actualizarse conforme aparezcan
actas civiles, parroquiales o documentos de archivo;
los árboles genealógicos en línea sirven como pistas,
no como prueba final.

\end{document}
